\label{chap:83-08}

% Unidad U012. Biografía estructural completa de pi.

\section{Construcción del carácter de clausura y trayectorias genealógicas estructural de \(\pi\)}
\label{p12-83-c8-s1:chap:biografia-pi-autonoma}

\subsection{Separación causal entre semilla, carácter y evaluación}

La construcción se divide en tres estratos:
\[
 \text{selección finita}
 \longrightarrow
 \text{construcción del carácter de clausura}
 \longrightarrow
 \text{evaluación arquimediana}.
\]
La expansión decimal no interviene en el primer estrato. La cadena de tipos
es
\begin{equation}
 104\,976
 \longrightarrow468\,\mathcal U_6
 \xrightarrow{\Psi}243\,\mathcal W_6
 \lhook\joinrel\longrightarrow\mathbb F_3^6,
 \qquad|\mathbb F_3^6|=729.
 \label{p12-83-c8-s1:eq:cadena-tipica-pi-autonoma}
\end{equation}
Los cardinales cuentan, respectivamente, condiciones iniciales
enriquecidas, emisiones distintas, palabras ternarias visibles y elementos
de la fibra ambiente completa.

Para \(U_6=(u_1,\ldots,u_6)\), la proyección es
\[
 \Psi(U_6)=((-u_1)\bmod3,\ldots,(-u_6)\bmod3).
\]
El censo de las \(43\) órbitas de \(D_3\) demuestra que la órbita de
clausura es la única cuyas seis orientaciones poseen a la vez:
\[
 \text{cinco elevaciones},\qquad
 \text{masa }1008,\qquad
 1008=432+4\cdot144.
\]
El calibre
\[
 \operatorname{diag}_c=6,
 \qquad
 \operatorname{card}=ES
\]
selecciona un único representante:
\begin{equation}
 \boxed{w_6^\pi=010211.}
 \label{p12-83-c8-s1:eq:semilla-pi-autonoma}
\end{equation}
La selección usa catálogo, multiplicidades, acción orbital y orientación;
no consulta un prefijo de \(\pi\).

\subsection{Elevaciones \(L_0,L_1\) y 1.ª frontera coinductiva}

Sobre vectores fila de \(\mathbb F_3^6\), sean
\[
L_0=
\begin{pmatrix}
2&2&2&1&2&1\\
2&1&2&2&1&1\\
1&0&0&1&0&2\\
0&0&1&1&1&0\\
2&2&2&0&2&1\\
2&1&2&1&0&0
\end{pmatrix},
\quad
L_1=
\begin{pmatrix}
2&1&2&0&2&2\\
1&2&1&1&2&0\\
1&2&1&1&1&2\\
0&1&0&0&2&1\\
2&0&2&1&0&1\\
0&2&2&2&1&1
\end{pmatrix}.
\]
Si \(b_1=w_6^\pi\), la recurrencia de bloques es
\[
 b_{k+1}=b_kL_{\varepsilon_k},
 \qquad
 (\varepsilon_1,\varepsilon_2,\varepsilon_3,\varepsilon_4)=(0,0,1,1),
\]
y la palabra acumulada se forma por concatenación:
\[
 w_{6(k+1)}=w_{6k}\mathbin{\|}b_{k+1}.
\]
No se multiplica una palabra de longitud \(6k\) por una matriz \(6\times6\).

Para el canal de clausura:
\[
\begin{array}{c|c}
\text{profundidad}&\text{bloque nuevo}\\ \hline
6&010211\\
12&012222\\
18&010211\\
24&002111\\
30&110221
\end{array}
\]
y
\begin{equation}
 \boxed{
 w_{30}^{\pi}
 =010211|012222|010211|002111|110221.}
 \label{p12-83-c8-s1:eq:w30-pi-autonomo}
\end{equation}
La frontera siguiente es
\[
 u_5^\pi=222220.
\]
Junto con los otros canales,
\[
 R_{36}=
 \begin{pmatrix}
 222220\\021222\\102011
 \end{pmatrix}.
\]
El símbolo \(R_{36}\) registra el sexto bloque de tres historias y abre una
frontera; no es un terminal.

La rama certificada de veinte bloques comienza
\begin{align*}
\pi:\;&010211|012222|010211|002111|110221|222220|111201|\\
&212121|200121|100100|101222|022212|012012|111210|\\
&121011|200220|120210|000101|022010|020111.
\end{align*}
Por ello, ni \(w_{30}\) ni \(R_{36}\) terminan la genealogía.

\subsection{Descomposición orientada de la microfibra de \texorpdfstring{\(\pi\)}{pi}: región central, \(4\) regiones laterales y multiplicidad \texorpdfstring{\(432+4\cdot144=1008\)}{432+4x144=1008}}

La microfibra
\[
 \mathscr R_\pi^{\rm micro}
 :=\{U\in\mathcal U_6:\Psi(U)=010211\}
\]
contiene exactamente cinco regiones:
\[
\begin{array}{c|c|c|c}
U_6&E_{\rm sig}&\operatorname{mult}&\text{papel}\\ \hline
501|614|498|169|272|272&555555&432&\perp\\
810|923|870|169|272|272&339555&144&++\\
810|983|810|169|272|272&393555&144&++\\
870|923|810|169|272|272&933555&144&++\\
870|923|810|418|674|521&933717&144&--
\end{array}
\]
En consecuencia,
\begin{equation}
 \boxed{
 5=1_\perp+3_{++}+1_{--},
 \qquad
 1008=432+4\cdot144.}
 \label{p12-83-c8-s1:eq:descomposiciones-cinco-regiones-pi}
\end{equation}
La primera identidad clasifica funciones de orientación. La segunda cuenta
multiplicidades de procedencia. No cuentan cinco copias indistinguibles.

\begin{figure}[htbp]
\centering
\resizebox{\textwidth}{!}{\input{figures/fig_cinco_regiones_pi}}
\caption{Microfibra de clausura de \(\pi\): región transversal,
\(3\) regiones de orientación positiva y retorno mutado. Las \(5\)
procedencias comparten la misma palabra cociente sin perder su genealogía.}
\label{fig:cinco-regiones-pi-microfibra-canonica}
\end{figure}

La región transversal es
\[
 501|614|498|169|272|272,
 \qquad E_{\rm sig}=555555,
 \qquad\operatorname{mult}=432.
\]
El retorno mutado es
\[
 870|923|810|418|674|521,
 \qquad E_{\rm sig}=933717,
 \qquad\operatorname{mult}=144.
\]
Intercambiarlos conservaría el censo \(3/1/1\), pero falsaría la estructura
regional completa.

\subsection{Descenso de un bloque nonádico completo al multigrafo de transferencia de fase}

Para \(U_6=(U_1,\ldots,U_6)\), sea
\[
 A(U_6)=(U_1,U_2,U_3),
 \qquad
 B(U_6)=(U_4,U_5,U_6).
\]
Cada emisión define una arista entre \(A(U_6)\) y \(B(U_6)\). El multigrafo
bruto tiene \(96\) vértices, distribuidos en componentes de tamaños
\[
 28,28,28,4,4,4.
\]
Se cocienta la fase interna mediante
\[
 (a,b,c)\sim(b,c,a)\sim(c,a,b).
\]
El representante es la rotación lexicográficamente mínima. El grafo
cocientado posee \(32\) vértices y dos componentes de tamaños \(28\) y \(4\).

Los anclajes son
\[
\begin{aligned}
U_{\rm APP}&=903|850|850|252|109|252,\\
U_e&=850|963|890|149|252|212,\\
U_\varphi&=830|943|830|109|252|252.
\end{aligned}
\]
\Needspace{9\baselineskip}
Numeramos las regiones
\[
\begin{aligned}
U_\pi^{(1)}&=810|923|870|169|272|272,\\
U_\pi^{(2)}&=810|983|810|169|272|272,\\
U_\pi^{(3)}&=870|923|810|169|272|272,\\
U_\pi^{(4)}&=870|923|810|418|674|521,\\
U_\pi^{(5)}&=501|614|498|169|272|272.
\end{aligned}
\]

\subsection{Los \(5\) ciclos orientados de clausura}

\subsubsection{Cierre de \texorpdfstring{\(U_\pi^{(1)}\)}{U-pi 1}}
\begin{align*}
&252|109|252|903|850|850\to
903|850|850|252|109|252\to\\
&252|109|252|923|870|810\to
810|923|870|169|272|272\to\\
&272|169|272|963|890|850\to
850|963|890|149|252|212\to\\
&212|149|252|943|830|830\to
830|943|830|109|252|252\to
252|109|252|903|850|850.
\end{align*}

\subsubsection{Cierre de \texorpdfstring{\(U_\pi^{(2)}\)}{U-pi 2}}
\begin{align*}
&903|850|850|252|109|252\to
272|169|272|903|850|850\to\\
&272|169|272|983|810|810\to
810|983|810|169|272|272\to\\
&272|169|272|963|890|850\to
850|963|890|149|252|212\to\\
&212|149|252|943|830|830\to
830|943|830|109|252|252\to
903|850|850|252|109|252.
\end{align*}

\subsubsection{Cierre de \texorpdfstring{\(U_\pi^{(3)}\)}{U-pi 3}}
\begin{align*}
&903|850|850|252|109|252\to
292|189|232|903|850|850\to\\
&292|189|232|923|810|870\to
870|923|810|169|272|272\to\\
&272|169|272|963|890|850\to
850|963|890|149|252|212\to\\
&212|149|252|943|830|830\to
830|943|830|109|252|252\to
903|850|850|252|109|252.
\end{align*}

\subsubsection{Cierre de \texorpdfstring{\(U_\pi^{(4)}\)}{U-pi 4}}
\begin{align*}
&903|850|850|252|109|252\to
272|169|272|903|850|850\to\\
&272|169|272|963|890|850\to
850|963|890|149|252|212\to\\
&212|149|252|923|810|870\to
870|923|810|418|674|521\to\\
&943|830|830|521|418|674\to
830|943|830|109|252|252\to
903|850|850|252|109|252.
\end{align*}

\subsubsection{Cierre de \texorpdfstring{\(U_\pi^{(5)}\)}{U-pi 5}}
\begin{align*}
&252|109|252|903|850|850\to
903|850|850|252|109|252\to\\
&614|498|501|252|109|252\to
501|614|498|169|272|272\to\\
&272|169|272|963|890|850\to
850|963|890|149|252|212\to\\
&212|149|252|943|830|830\to
830|943|830|109|252|252\to
252|109|252|903|850|850.
\end{align*}

\begin{theorem}[Minimalidad de los cinco cierres]
\label{p12-83-c8-s1:thm:minimalidad-cierres-pi-autonoma}
Para cada \(i\in\{1,\ldots,5\}\), la longitud mínima de un paseo cerrado
que recorre simultáneamente
\[
 U_\pi^{(i)},\quad U_e,\quad U_\varphi,\quad U_{\rm APP}
\]
en el multigrafo cocientado es ocho.
\end{theorem}

\begin{proof}
Fijado \(i\), sea
\[
 E_*=\{U_\pi^{(i)},U_e,U_\varphi,U_{\rm APP}\}.
\]
Se considera el espacio finito
\[
 \mathscr S_i=\{(v,M):v\in V(G),\ M\subseteq E_*\}.
\]
El primer componente registra el vértice y el segundo las aristas objetivo
ya recorridas. Si una arista \(e\) lleva \(v\) a \(v'\), la transición es
\[
 (v,M)\longmapsto
 \bigl(v',M\cup(E_*\cap\{e\})\bigr).
\]
La búsqueda en anchura enumera longitudes en orden creciente. El primer
retorno a \((v,E_*)\) aparece en longitud ocho para las cinco regiones.
Los paseos impresos realizan la cota; la optimalidad de la búsqueda excluye
una longitud menor.
\end{proof}

La función de clausura queda codificada por
\[
 \mathfrak C_\pi=
 \bigl(\mathscr R_\pi^{\rm micro},
 w_\pi,\rho_\pi,m_\pi,\ell_{\rm cl}\bigr),
\]
\[
 \rho_\pi=(\perp,++ ,++ ,++ ,--),
 \qquad
 m_\pi=(432,144,144,144,144),
 \qquad
 \ell_{\rm cl}=8.
\]

\subsection{Operador de incidencia y media vuelta orientada}

La sombra de incidencia es
\[
 A_W=
 \begin{pmatrix}
 0&1&1&1&1&1\\
 1&0&1&2&2&1\\
 1&1&0&1&2&2\\
 1&2&1&0&1&2\\
 1&2&2&1&0&1\\
 1&1&2&2&1&0
 \end{pmatrix}.
\]
La multiplicación en \(\mathbb F_3\) da
\begin{equation}
 \boxed{A_W^2=2I_6=-I_6.}
 \label{p12-83-c8-s1:eq:aw-cuadrado-menos-identidad-pi}
\end{equation}
Los productos entre filas y columnas distintas se anulan módulo tres,
mientras cada producto diagonal vale dos. Una aplicación realiza un cuarto
de giro ternario; dos aplicaciones realizan la inversión orientada.

\subsection{La pendiente 1/5 forzada por las \(5\) regiones}

Sea
\[
 \lambda=(\lambda_1,\ldots,\lambda_5)^{\mathsf T}
 \in\mathbb Q_{\geq0}^5,
 \qquad
 \mathbf1^{\mathsf T}\lambda=1.
\]
Olvidar la etiqueta regional aplica el simetrizador
\[
 P_5=\frac15\mathbf1\mathbf1^{\mathsf T}.
\]
Para todo vector normalizado,
\[
 P_5\lambda=\frac15\mathbf1.
\]
La invariancia \(P_5\lambda=\lambda\) tiene, por tanto, una única solución:
\begin{equation}
 \boxed{\lambda=\frac15\mathbf1,\qquad t_1=\frac15.}
 \label{p12-83-c8-s1:eq:pendiente-primitiva-pi}
\end{equation}
El entero cinco procede de la microfibra regional; no se reconoce en una
fórmula de \(\pi\).

\subsection{Composición racional y compensador \texorpdfstring{\(239\)}{239}}

La ley de composición de pendientes es
\[
 t\star s:=\frac{t+s}{1-ts}.
\]
La primera duplicación da
\[
 t_2=t_1\star t_1
 =\frac{2/5}{1-1/25}
 =\frac5{12}.
\]
La segunda da
\[
 t_4=t_2\star t_2
 =\frac{10/12}{1-25/144}
 =\frac{120}{119}.
\]
El retorno a la diagonal unitaria exige \(q>0\) tal que
\[
 \frac{120/119-1/q}{1+(120/119)/q}=1.
\]
Multiplicar por \(119q\) produce
\[
 120q-119=119q+120,
\]
y fuerza
\begin{equation}
 \boxed{q=239.}
 \label{p12-83-c8-s1:eq:compensador-239-pi}
\end{equation}

\subsection{Identidad gaussiana y 1.ª media vuelta}

Definimos
\begin{equation}
 \Theta_\pi
 :=16\arctan\frac15-4\arctan\frac1{239}.
 \label{p12-83-c8-s1:eq:theta-pi-autonoma}
\end{equation}
Los cuadrados sucesivos son
\[
\begin{array}{c|r@{}l}
\text{potencia}&\multicolumn{2}{c}{\text{entero gaussiano}}\\ \hline
(5+i)^2&24&+10i\\
(5+i)^4&476&+480i\\
(5+i)^8&-3824&+456960i\\
(5+i)^{16}&-208797818624&-3494830080i\\
(239-i)^2&57120&-478i\\
(239-i)^4&3262465916&-54606720i
\end{array}
\]
y la multiplicación final da
\begin{equation}
 \boxed{
 (5+i)^{16}(239-i)^4
 =-681386607803576157184.}
 \label{p12-83-c8-s1:eq:identidad-gaussiana-pi-autonoma}
\end{equation}
La parte imaginaria es cero y la real es negativa. Por tanto,
\[
 \Theta_\pi\equiv\pi\pmod{2\pi}.
\]
Para \(0<x<1\),
\[
 x-\frac{x^3}{3}<\arctan x<x.
\]
De aquí
\[
 \Theta_\pi>
 16\left(\frac15-\frac1{3\cdot5^3}\right)-\frac4{239}>3
\]
y
\[
 \Theta_\pi<\frac{16}{5}<4.
\]
La única orientación negativa en ese intervalo es la primera media vuelta
positiva:
\begin{equation}
 \boxed{\Theta_\pi=\pi.}
 \label{p12-83-c8-s1:eq:reconocimiento-pi-exacto}
\end{equation}

El orden causal es
\[
 \text{cinco regiones}
 \to\frac15
 \to\frac5{12}
 \to\frac{120}{119}
 \to239
 \to\Theta_\pi
 \to\pi.
\]

\subsection{Horquillas racionales de precisión arbitraria}

Para \(0<x<1\), sea
\[
 S_n(x)=\sum_{r=0}^{n}(-1)^r\frac{x^{2r+1}}{2r+1}.
\]
El criterio de Leibniz da
\[
 S_{2m+1}(x)<\arctan x<S_{2m}(x).
\]
Definimos
\[
 a_m^-=S_{2m+1}(1/5),\quad a_m^+=S_{2m}(1/5),
\]
\[
 b_m^-=S_{2m+1}(1/239),\quad b_m^+=S_{2m}(1/239),
\]
y
\begin{equation}
 \ell_{\pi,m}=16a_m^- -4b_m^+,
 \qquad
 h_{\pi,m}=16a_m^+ -4b_m^-.
 \label{p12-83-c8-s1:eq:horquilla-pi-autonoma}
\end{equation}
Entonces
\[
 \ell_{\pi,m}<\pi<h_{\pi,m}
\]
y
\begin{equation}
 h_{\pi,m}-\ell_{\pi,m}
 =
 \frac{16}{(4m+3)5^{4m+3}}
 +\frac{4}{(4m+3)239^{4m+3}}
 \longrightarrow0.
 \label{p12-83-c8-s1:eq:anchura-horquilla-pi}
\end{equation}
Las cotas inferiores crecen y las superiores decrecen.

Como \(3<\pi<4\), el lector ternario actúa sobre
\[
 r_\pi:=\pi-3\in(0,1),
\]
con horquilla
\[
 [\ell_{\pi,m}-3,h_{\pi,m}-3].
\]

\subsection{Producción por \texorpdfstring{$\operatorname{Ext}$}{Ext} y certificación racional posterior en la fibra ambiente}

Supóngase construido por \(\operatorname{Ext}\), mediante el transporte
graduado del TPK sobre el estado enriquecido, un prefijo ternario de longitud
\(L_K=6K\), representado por \(N_K\). La misma ley de prolongación produce
primero la extensión superviviente única \(N_{K+1}\), conservando hoja,
ruta, residuo, cociente, acarreo, frontera y memoria. Los \(729\)
subcilindros de la fibra ambiente en la que se publica esa salida son
\[
 I_{K,u}=
 \left[
 \frac{729N_K+u}{3^{L_K+6}},
 \frac{729N_K+u+1}{3^{L_K+6}}
 \right),
 \qquad0\leq u<729.
\]
Una vez producido \(N_{K+1}\), el orden \(m=m(K)\) se aumenta hasta que la
horquilla racional quede contenida en un único subcilindro y certifique así,
a posteriori, la coordenada ya publicada por \(\operatorname{Ext}\). Sean
\[
 \ell_{\pi,K}:=\ell_{\pi,m(K)}-3,
 \qquad
 h_{\pi,K}:=h_{\pi,m(K)}-3.
\]
Los extremos enteros compatibles son
\[
 j_{\min}=
 \max\!\left(
 729N_K,\left\lfloor
 \ell_{\pi,K}3^{L_K+6}\right\rfloor
 \right),
\]
\[
 j_{\max}=
 \min\!\left(
 729N_K+728,\left\lceil
 h_{\pi,K}3^{L_K+6}\right\rceil-1
 \right).
\]
La condición de certificación unitaria es
\begin{equation}
 \boxed{j_{\min}=j_{\max},\qquad N_{K+1}=j_{\min}.}
 \label{p12-83-c8-s1:eq:seleccion-unitaria-subcilindro-pi}
\end{equation}
La igualdad expresa el barrido exhaustivo de la fibra ambiente y reconoce el
subcilindro numérico que contiene la extensión previamente producida. En
particular, \(N_{K+1}=j_{\min}\) se verifica aquí como identidad de
certificación: no define \(N_{K+1}\), no selecciona su bloque y no alimenta el
generador. La existencia y la unicidad de la extensión pertenecen al
transporte \(\operatorname{Ext}\) del estado enriquecido; la horquilla y la
fibra de \(729\) elementos certifican posteriormente su localización
arquimediana.

\subsection{Prolongación coinductiva de la sección de \(\pi\) mediante el sistema inverso de historias supervivientes}

Sea \(\mathcal H_m^\pi\) el conjunto de historias enriquecidas a profundidad
\(m\), y sea
\[
 p_{n,m}:\mathcal H_n^\pi\longrightarrow\mathcal H_m^\pi
\]
el truncamiento. Definimos
\[
 \operatorname{Surv}_m^{(N),\pi}
 :=p_{N,m}(\mathcal H_N^\pi),
\]
\[
 \operatorname{Surv}_m^\pi
 :=\bigcap_{N\geq m}\operatorname{Surv}_m^{(N),\pi}.
\]
Cuando la relación \(\operatorname{Ext}\) es no vacía, univaluada y
compatible en todos los niveles, los conjuntos son unitarios y satisfacen
\[
 p_{m+1,m}(\operatorname{Surv}_{m+1}^{\pi})
 =\operatorname{Surv}_{m}^{\pi}.
\]
En ese dominio de prolongabilidad existe, por tanto, una única sección
\[
 x_\pi\in
 X_\pi:=\varprojlim_m\operatorname{Surv}_m^\pi.
\]
Los cilindros están encajados y sus diámetros son \(3^{-6K}\), que tiende a
cero. Su intersección contiene un único punto, identificado por el carácter
angular con \(r_\pi=\pi-3\).

Para toda profundidad decimal \(M\), puede elegirse un nivel \(K(M)\) cuyo
cilindro queda dentro de una sola celda decimal de diámetro \(10^{-M}\).
Los prefijos profundos truncan exactamente a los anteriores.

\subsection{Caracterización del canal coinductivo de \(\pi\): selección de \(w_6^\pi\), \(5\) regiones, inversión orientada y unicidad de la sección límite}

\begin{theorem}[Caracterización del canal de clausura]
\label{p12-83-c8-s1:thm:caracterizacion-final-pi-autonoma}
Partiendo del catálogo, la acción \(D_3\), el calibre orientado, los
operadores \(L_0,L_1,A_W\), la ley de pendientes y la filtración racional:
\begin{enumerate}
 \item se selecciona \(w_6^\pi=010211\);
 \item se construyen cinco regiones con las identidades de
       \eqref{p12-83-c8-s1:eq:descomposiciones-cinco-regiones-pi};
 \item cada región pertenece a un cierre mínimo de longitud ocho;
 \item \(A_W^2=-I_6\) fija la inversión orientada;
 \item la simetrización regional fuerza \(1/5\);
 \item la composición fuerza \(120/119\);
 \item el retorno unitario fuerza \(239\);
 \item la identidad gaussiana reconoce la primera media vuelta;
 \item \(\operatorname{Ext}\), mediante el transporte graduado del TPK,
       produce la sección compatible a profundidad arbitraria, mientras las
       horquillas y la fibra de \(729\) elementos certifican después, en cada
       profundidad finita, el bloque ya producido.
\end{enumerate}
Su evaluación es
\[
 \boxed{\pi=16\arctan\frac15-4\arctan\frac1{239}.}
\]
\end{theorem}

La cadena completa es
\[
 \mathrm{APP}\to\mathrm{TRIT}\to\mathrm{TPK}
 \to104\,976\to468\to243\to w_6^\pi
 \to w_{30}^\pi\to R_{36}\to\Gamma_9
 \to X_\pi\to\pi.
\]

\begin{criterion}[Falsación del canal de clausura]
La derivación queda refutada si:
\begin{enumerate}
 \item la órbita de clausura deja de ser única;
 \item desaparece una de las cinco regiones o cambia su identidad completa;
 \item uno de los cinco cierres tiene longitud menor que ocho;
 \item \(A_W^2\ne2I_6\);
 \item el simetrizador admite otra distribución normalizada;
 \item las composiciones no producen \(5/12\), \(120/119\) y \(239\);
 \item falla la identidad gaussiana;
 \item una horquilla deja de contener el carácter angular;
 \item un nivel no recorre los \(729\) elementos de la fibra ambiente o no
       localiza exactamente una extensión;
 \item una expansión objetivo interviene antes del reconocimiento.
\end{enumerate}
\end{criterion}
